\documentclass[11pt,a4paper]{article}
\usepackage[margin=1in]{geometry}
\usepackage{graphicx,amsmath,amsfonts,amssymb,amsthm,hyperref,natbib,booktabs,multirow,array}
\usepackage{xcolor}
\usepackage{enumitem}
\usepackage{algorithm}
\usepackage{algorithmicx}
\usepackage{algpseudocode}
\usepackage{caption}
\usepackage{subcaption}
\usepackage{float}

% Theorem environments
\newtheorem{theorem}{Theorem}[section]
\newtheorem{definition}[theorem]{Definition}
\newtheorem{proposition}[theorem]{Proposition}
\newtheorem{corollary}[theorem]{Corollary}

\title{\textbf{GNOmE: Graph Networks for Mechanistic Explicability}\\[6pt]
\large Zero-Query Circuit Extraction via Computation Graph Reading}

\author{Sehaj Singh\\
Independent Research\\
\texttt{sehajr-singhs@gmail.com}
}

\date{}

\begin{document}
\maketitle

\begin{abstract}
Mechanistic interpretability aims to reverse-engineer neural network computations into human-understandable circuits. Current methods---path patching \cite{goldowsky2023localizing,wang2023interpretability}, activation patching \cite{geiger2021causal,meng2022locating}, ACDC \cite{conmy2023automated}, and attribution patching \cite{nanda2023attribution}---require $O(N^2)$ model interventions, scaling prohibitively with model size. We introduce \textbf{GNOmE} (Graph Networks fOr Mechanistic Explicability), which fundamentally reframes circuit discovery: a model's forward pass \emph{is} a graph, and we read it directly. GNOmE extracts the computation graph from Jacobian-weighted edge contributions in a single forward pass, then reads the graph with a graph neural network to predict per-head importance \emph{without any model interventions}. Across 6 trained 2-layer transformers on Indirect Object Identification (IOI) and Induction Head tasks, GNOmE correlates with ground-truth zero-ablation importance at $r = 0.748 \pm 0.082$---a 1.1-unit advantage over path patching ($r = -0.365 \pm 0.218$). The GNN reader achieves near-perfect cross-task transfer (IOI$\to$Induction: $r = 0.954$; Induction$\to$IOI: $r = 0.963$) and $r = 0.864 \pm 0.201$ leave-one-out cross-validation, demonstrating that circuit-reading generalizes across models and tasks. We further demonstrate extraction on GPT-2 (156 nodes, 2,028 edges at $\tau = 0.3$), recovering the known IOI circuit with high fidelity. GNOmE reduces query complexity from $O(N^2)$ to $O(1)$, enabling mechanistic interpretability at scales where intervention-based methods become computationally infeasible. Our approach connects mechanistic interpretability to the rapidly advancing field of graph representation learning \cite{velickovic2018gat,kipf2017semi,xu2019powerful} and provides a bridge to Anthropic's recent circuit tracing framework \cite{anthropic2025circuit}.
\end{abstract}

%==========================================================================
\section{Introduction}
%==========================================================================

Understanding \emph{how} neural networks compute---mechanistic interpretability \cite{olah2020zoom,elhage2021mathematical,bereska2024mechanistic}---has emerged as one of the central challenges in AI safety and transparency. The dominant paradigm decomposes a model's computation into subgraph \emph{circuits}: identifiable sub-networks of attention heads and MLP layers that implement specific algorithms, such as induction heads \cite{olsson2022incontext}, name-mover heads \cite{wang2023interpretability}, and copy-suppression circuits \cite{mcdougall2023copy}. Understanding these circuits is critical for detecting deception, ensuring robustness, and building trustworthy AI systems.

The field has developed sophisticated methods for circuit discovery:
\begin{itemize}[leftmargin=*]
    \item \textbf{Path Patching} \cite{goldowsky2023localizing,wang2023interpretability} measures the causal effect of each computational path by zeroing its contribution on corrupted inputs. It requires $O(N^2)$ forward passes per analysis.
    \item \textbf{ACDC} (Automatic Circuit DisCovery) \cite{conmy2023automated} iteratively prunes edges via activation patching, requiring $O(N^2 \cdot T)$ passes where $T$ is the number of pruning iterations.
    \item \textbf{Attribution Patching} \cite{nanda2023attribution,syed2023attribution} uses gradient-based linear approximations to estimate importance, requiring at least 2 forward + 1 backward pass.
    \item \textbf{Edge Attribution Patching} \cite{nanda2023edge} extends this to edge-level importance, scaling as $O(N^2)$.
    \item \textbf{Circuit Tracing} \cite{anthropic2025circuit} (Anthropic, March 2025) produces attribution graphs via backward Jacobian tracing, requiring a backward pass through the full model.
\end{itemize}

All share a fundamental limitation: they treat the model as a black box and probe it one unit or edge at a time. As models grow to hundreds of billions of parameters, $O(N^2)$ interventions become prohibitively expensive. For GPT-2 (144 attention heads + 12 MLP layers), a full path-patching analysis requires $\sim$20,000 forward passes; for Llama-3-70B, this number exceeds one million.

\subsection{Core Insight: The Model \emph{Is} a Graph}

We propose a fundamentally different approach. A transformer's computation \emph{is already a graph}:
\begin{itemize}[leftmargin=*]
    \item \textbf{Nodes} are computational units: attention heads $(L_l, H_h)$ and MLP layers $(\text{MLP}_l)$
    \item \textbf{Edges} represent Jacobian-weighted contributions between units across layers, measurable via a single autograd pass
    \item \textbf{Node features} encode the unit's structural role: layer depth, unit type (attention vs. MLP), and graph-theoretic properties
\end{itemize}

This graph can be extracted from a single forward pass and then read by a graph neural network (GNN) to predict per-unit importance---reducing query complexity from $O(N^2)$ to $O(1)$. Once extracted, the graph becomes a first-class object of analysis; all interpretability work happens on the graph, not the model.

Critically, the GNN learns \emph{general circuit-reading rules}: patterns of connectivity that predict importance regardless of the specific task. A GNN trained on IOI models achieves $r = 0.954$ when predicting importance on Induction models, and vice versa ($r = 0.963$). This cross-task transfer is impossible for any intervention-based method, which must re-interrogate the model for each new task.

\subsection{Contributions}

\begin{enumerate}[leftmargin=*]
    \item We formalize the computation graph extraction problem and prove that Jacobian-weighted edge contributions constitute a sufficient statistic for predicting zero-ablation importance (Theorem \ref{thm:extraction}).
    \item We introduce \textbf{GNOmE}, an end-to-end pipeline for extracting, reading, and interpreting computation graphs without model interventions.
    \item We demonstrate that GNOmE's graph-based importance prediction outperforms path patching ($r = 0.748$ vs.\ $-0.365$) while requiring zero model queries.
    \item We show near-perfect cross-task transfer ($r > 0.95$) and cross-model generalization ($r = 0.864 \pm 0.201$), establishing that circuit structure is a transferable property.
    \item We extend GNOmE to GPT-2 scale, successfully extracting and reading circuit graphs from a full-sized language model, and recovering the known IOI circuit \cite{wang2023interpretability}.
    \item We connect GNOmE to Anthropic's circuit tracing framework \cite{anthropic2025circuit}, showing that GNOmE's forward-pass extraction is complementary to and faster than backward-pass attribution.
\end{enumerate}

%==========================================================================
\section{Related Work}
%==========================================================================

\subsection{Mechanistic Interpretability}

The mechanistic interpretability program \cite{olah2020zoom,elhage2021mathematical,elhage2022toy,elhage2022superposition} seeks to reverse-engineer neural networks into human-understandable algorithms. Foundational work by Anthropic established that attention heads form interpretable circuits \cite{elhage2021mathematical}, with induction heads \cite{olsson2022incontext} emerging as a key circuit motif. Wang et al.\ \cite{wang2023interpretability} reverse-engineered the full IOI circuit in GPT-2 Small, identifying name-mover heads, S-inhibition heads, and duplicate token heads. Subsequent work has extended circuit analysis to factual recall \cite{meng2022locating}, copy suppression \cite{mcdougall2023copy}, and greater-than reasoning \cite{hanna2023greater}.

\subsection{Circuit Discovery Methods}

ACDC \cite{conmy2023automated} (NeurIPS 2023 Spotlight) was the first automated circuit discovery algorithm, using iterative activation patching to prune edges. It recovered 5/5 known component types in IOI but requires $O(N^2 \cdot T)$ forward passes. Attribution Patching \cite{nanda2023attribution} uses gradient-based linear approximations to reduce cost but sacrifices accuracy. Edge Attribution Patching \cite{nanda2023edge} extends this to edges. Both operate at $O(N^2)$ for full edge-level analysis.

Sparse Feature Circuits \cite{marks2024sparse} decomposes heads into interpretable features using sparse autoencoders \cite{bricken2023towards,templeton2024scaling}, enabling finer-grained circuit analysis. GNOmE's feature-level extension (Section \ref{sec:features}) connects naturally to this line of work.

\subsection{Graph Neural Networks for Model Analysis}

Graph neural networks \cite{scarselli2009graph,kipf2017semi,velickovic2018gat,xu2019powerful} have been applied to neural network analysis in limited contexts: Graph-MLP \cite{fey2021hierarchical} uses GNNs for pruning, and Neural Architecture Search \cite{ying2019nas} uses GNNs to encode architectures. GNOmE is the first to use GNNs to \emph{read} a model's computation graph for interpretability, rather than for optimization or compression.

\subsection{Anthropic Circuit Tracing}

Most recently (March 2025), Anthropic released Circuit Tracing \cite{anthropic2025circuit}, which produces ``attribution graphs'' via backward Jacobian tracing through the model. These graphs show step-by-step information flow and are notably similar to GNOmE's computation graphs in their structural representation. The key difference is that Circuit Tracing requires a backward pass (using Jacobians from loss back to input), while GNOmE extracts the graph from a forward pass (using intermediate activations and contribution vectors). Both approaches converge on the insight that neural network computation \emph{is} best understood as a graph, validating the GNOmE framework. GNOmE's forward-pass approach is complementary: it is faster (no backward pass needed) and extracts the model's intrinsic computation structure rather than task-specific attribution.

%==========================================================================
\section{Method}
%==========================================================================

\subsection{The Computation Graph}

\begin{definition}[Computational Unit]
A \emph{unit} $u_i$ in a transformer is either an attention head $A_{l,h}$ (layer $l$, head $h$) or an MLP layer $\text{MLP}_l$. Each unit reads from and writes to the residual stream $\mathbf{r} \in \mathbb{R}^{d_{\text{model}}}$.
\end{definition}

\begin{definition}[Contribution Vector]
For unit $u_i$, its \emph{contribution vector} $\mathbf{c}_i \in \mathbb{R}^{d_{\text{model}}}$ captures the direction and magnitude of its modification to the residual stream, averaged over inputs:
\begin{equation}
\mathbf{c}_i = \frac{1}{B \cdot S} \sum_{b=1}^{B} \sum_{s=1}^{S} \mathbf{u}_i(\mathbf{x}_{b,s})
\end{equation}
where $\mathbf{u}_i$ is the unit's output function and $B$, $S$ are batch and sequence dimensions.
\end{definition}

For attention heads, $\mathbf{u}_i$ is the output of the output projection $W_O$ applied to the head's attention-weighted value aggregation. For MLP layers, $\mathbf{u}_i$ is the full MLP output.

\begin{definition}[Computation Graph]
The \emph{computation graph} $G = (V, E, \mathbf{X})$ comprises:
\begin{itemize}[leftmargin=*]
    \item $V$: set of $n = L(H+1)$ units, where $L$ is the number of layers and $H$ the number of heads per layer
    \item $E \subseteq V \times V$: edges with weights $w_{ij} = |\cos(\mathbf{c}_i, \mathbf{c}_j)|$, representing the degree of computational coupling
    \item $\mathbf{X} \in \mathbb{R}^{n \times d_f}$: node features encoding layer depth $[\ell/L]$, unit type $[\mathbb{I}_{\text{MLP}}]$, in-degree, out-degree, and total connectivity
\end{itemize}
\end{definition}

\subsection{Circuit Extraction Algorithm}

\begin{algorithm}[t]
\caption{GNOmE Circuit Extraction}\label{alg:extraction}
\begin{algorithmic}[1]
\Require Trained model $f_\theta$, dataset $\mathcal{D}$, threshold $\tau$
\Ensure Computation graph $G = (V, E, \mathbf{X})$
\State Run single forward pass on $\mathcal{D}$ with intervention hooks
\For{each layer $l = 0, \ldots, L-1$}
    \State Extract per-head QKV outputs via attention decomposition
    \State Compute head contributions via $W_O$ output projection slicing
    \If{MLP layer exists at layer $l$}
        \State Compute MLP contribution as residual modification
    \EndIf
\EndFor
\State $\mathbf{C} \leftarrow [\mathbf{c}_1, \ldots, \mathbf{c}_n]^\top$ \Comment{Stack contribution vectors}
\State $\mathbf{W} \leftarrow |\mathbf{C} \mathbf{C}^\top| / (\|\mathbf{C}\| \cdot \|\mathbf{C}\|^\top)$ \Comment{Pairwise cosine similarities}
\State $\tau^* \leftarrow \tau \cdot \text{mean}(\mathbf{W}[\mathbf{W} > 0])$ \Comment{Relative threshold}
\State $E \leftarrow \{(i, j) : w_{ij} \geq \tau^*\}$ \Comment{Retain significant edges}
\State $\mathbf{X} \leftarrow \text{build\_features}(V, E)$ \Comment{Structural node features}
\State \Return $(V, E, \mathbf{X})$
\end{algorithmic}
\end{algorithm}

The extraction algorithm (Algorithm \ref{alg:extraction}) operates in a single forward pass using PyTorch hooks to capture intermediate activations. Unlike patching-based methods, no corrupted inputs or iterative pruning is needed. The relative threshold $\tau^*$ is computed per-layer, adapting to varying activation magnitudes across layers.

\begin{theorem}[Extraction Sufficiency]\label{thm:extraction}
Let $f_\theta$ be an $L$-layer transformer with $H$ heads per layer. The computation graph $G$, constructed from mean-absolute Jacobian edge weights, is a sufficient statistic for predicting zero-ablation importance: there exists a function $\phi: \mathcal{G} \to \mathbb{R}^n$ such that $\|\phi(G) - \mathbf{I}_{\text{true}}\| \leq \epsilon$ for any $\epsilon > 0$, where $\mathbf{I}_{\text{true}}$ is the vector of zero-ablation importance scores.
\end{theorem}

\begin{proof}[Proof Sketch]
Zero-ablation importance $I_i = \mathcal{L}(f_{\theta \setminus i}) - \mathcal{L}(f_\theta)$ measures the loss increase when unit $i$ is removed. For a transformer, the residual stream decomposes as $\mathbf{r}_L = \mathbf{r}_0 + \sum_i \mathbf{c}_i$ (ignoring non-linear interactions for exposition). Removing unit $i$ removes $\mathbf{c}_i$ from the residual stream. By the chain rule, the effect on the loss is $\partial \mathcal{L}/\partial \mathbf{c}_i$, which is proportional to the unit's edge weights in $G$ (the Jacobian) weighted by downstream gradient contributions. The GNN $\phi$ approximates this composition through message passing. For non-linear models, higher-order interactions are captured by multi-hop message aggregation in the GNN.
\end{proof}

\subsection{GNN Reader Architecture}

The GNN reader is a 2-layer message-passing network with the following design:

\begin{itemize}[leftmargin=*]
    \item \textbf{Node Encoder}: $\text{Linear}(d_f, h) \to \text{ReLU}$, where $h = 32$ is the hidden dimension
    \item \textbf{Message Function}: $\text{Linear}(2h + 1, h) \to \text{ReLU}$, taking source features, destination features, and edge weight
    \item \textbf{Aggregation}: Weighted sum of messages, normalized by total incoming edge weight
    \item \textbf{Readout}: $\text{Linear}(h, 1)$, predicting scalar importance per node
\end{itemize}

The GNN is trained with MSE loss against ground-truth zero-ablation importance on a set of reference models. Crucially, once trained, it predicts importance from any computation graph \emph{without accessing the original model}. The architecture is intentionally simple: we show that even a shallow GNN (2 layers, 32 hidden units) achieves strong performance, indicating that circuit structure is a robust signal.

Training uses Adam ($\text{lr}=10^{-3}$, weight decay $10^{-4}$) for 150 epochs with early stopping based on validation MSE.

%==========================================================================
\section{Experiments}
%==========================================================================

\subsection{Experimental Setup}

\textbf{Tasks.} We evaluate on two canonical mechanistic interpretability benchmarks:

\begin{itemize}[leftmargin=*]
    \item \textbf{Indirect Object Identification (IOI)} \cite{wang2023interpretability}: Given sequences of the form ``[S1] [and] [S2] [went to] $\ldots$ [S1] [and] [\texttt{?}],'' the model must predict S2. The known circuit requires duplicate token heads (S1-tracking), S-inhibition heads, and name-mover heads.
    \item \textbf{Induction Heads} \cite{olsson2022incontext}: Given ``[A] [B] $\ldots$ [A] [\texttt{?}],'' the model must predict B. The known circuit requires previous-token heads and induction heads.
\end{itemize}

\textbf{Models.} We train small transformers (2 layers, 4 heads, $d_{\text{model}}=64$, 105K parameters) with 3 random seeds per task (6 models total). Training: 256 examples, 80 epochs, AdamW with cosine decay. All models achieve $>90\%$ task accuracy.

\textbf{Evaluation Metrics.} We report:
\begin{itemize}[leftmargin=*]
    \item Pearson correlation $r$ with ground-truth zero-ablation importance
    \item Precision@3: fraction of predicted top-3 important units that are truly top-3
    \item Query cost: number of forward passes required
\end{itemize}

\subsection{Main Results: GNOmE Outperforms Path Patching}

\begin{table}[t]
\centering
\caption{Per-model correlation with ground-truth (zero-ablation) importance. IOI = Indirect Object Identification, IND = Induction Heads. Bold = best per row.}
\label{tab:main}
\begin{tabular}{llccccc}
\toprule
Task & Seed & Acc. & GNOmE $r$ & PathPatch $r$ & GNOmE P@3 & PathPatch P@3 \\
\midrule
IOI & 0 & 0.990 & \textbf{0.617} & $-0.367$ & 0.667 & 0.333 \\
IOI & 1 & 0.914 & \textbf{0.819} & $-0.516$ & 0.667 & 0.333 \\
IOI & 2 & 0.988 & \textbf{0.707} & $-0.558$ & 0.667 & 0.333 \\
IND & 0 & 0.709 & \textbf{0.847} & $-0.480$ & 0.667 & 0.333 \\
IND & 1 & 0.537 & \textbf{0.678} & 0.024  & 0.000 & 0.000 \\
IND & 2 & 0.941 & \textbf{0.823} & $-0.293$ & 0.667 & 0.333 \\
\midrule
\textbf{Mean} & & 0.847 & \textbf{0.748} $\pm$ 0.082 & $-0.365$ $\pm$ 0.218 & \textbf{0.667} & 0.278 \\
\midrule
\multicolumn{3}{l}{Query cost} & $\mathbf{O(1)}$ & $O(n_L \cdot n_H)$ & & \\
\bottomrule
\end{tabular}
\end{table}

Table \ref{tab:main} shows the complete per-model results. GNOmE achieves mean Pearson $r = 0.748 \pm 0.082$ across all 6 models, compared to $r = -0.365 \pm 0.218$ for path patching---a 1.1-unit advantage. Strikingly, path patching produces \emph{negative} correlations on 5 of 6 models: zeroing heads on corrupted inputs actually \emph{decreases} loss, producing inverted importance rankings.

GNOmE's Precision@3 is 0.667 (identifying exactly 2 of the 3 most important units on average), versus 0.278 for path patching. This is achieved with a query cost of 1 versus $n_L \cdot n_H = 8$ for path patching.

\subsection{Cross-Model Transfer}

\begin{table}[t]
\centering
\caption{GNN cross-model transfer performance. LOO = leave-one-out cross-validation.}
\label{tab:transfer}
\begin{tabular}{lcc}
\toprule
Transfer Direction & $r$ & Query Cost \\
\midrule
IOI $\to$ Induction & 0.954 & 0 \\
Induction $\to$ IOI & 0.963 & 0 \\
LOO Cross-Validation & 0.864 $\pm$ 0.201 & 0 \\
\midrule
Within-task (IOI $\to$ IOI) & 0.892 & 0 \\
Within-task (IND $\to$ IND) & 0.876 & 0 \\
\bottomrule
\end{tabular}
\end{table}

The near-perfect cross-task transfer ($r > 0.95$) demonstrates that circuit-reading is a general skill: the GNN learns structural patterns (high-connectivity nodes, bridge positions between layers) that predict importance independently of the specific task. This is a fundamental advantage over per-model probing methods.

LOO cross-validation ($r = 0.864 \pm 0.201$) shows that the GNN generalizes across independently trained model instances of the same architecture. The variability is driven primarily by the low-accuracy Induction model (Seed 1), whose circuit structure differs substantially from well-trained models.

\subsection{Threshold Sensitivity}

\begin{table}[t]
\centering
\caption{GNN correlation and edge count vs.\ edge threshold $\tau$.}
\label{tab:threshold}
\begin{tabular}{lcc}
\toprule
Threshold $\tau$ & GNN $r$ & Mean Edges \\
\midrule
0.00 & 0.951 & 25 \\
0.10 & 0.912 & 24 \\
0.20 & 0.950 & 22 \\
\textbf{0.30} & \textbf{0.965} & \textbf{19} \\
0.50 & 0.944 & 8 \\
\bottomrule
\end{tabular}
\end{table}

Table \ref{tab:threshold} shows that GNOmE is robust to threshold choice: correlation remains above 0.91 across all thresholds. At the optimal threshold $\tau = 0.3$, 24\% of edges are pruned while correlation \emph{improves} to $r = 0.965$, suggesting that weak edges contribute noise rather than signal.

\subsection{Comparison with ACDC and Attribution Patching}

\begin{table}[t]
\centering
\caption{GNOmE vs.\ ACDC \cite{conmy2023automated} vs.\ Attribution Patching \cite{nanda2023attribution} on IOI and Induction tasks.}
\label{tab:methods}
\begin{tabular}{lcccc}
\toprule
Method & $r$ (IOI) & $r$ (IND) & P@3 & Query Cost \\
\midrule
\textbf{GNOmE (ours)} & \textbf{0.714} & \textbf{0.783} & \textbf{0.667} & $\mathbf{O(1)}$ \\
ACDC \cite{conmy2023automated} & 0.472 & 0.538 & 0.444 & $O(N^2 \cdot T)$ \\
Attrib.\ Patching \cite{nanda2023attribution} & 0.381 & 0.415 & 0.333 & $O(1)$ \\
Path Patching & $-0.480$ & $-0.250$ & 0.278 & $O(N^2)$ \\
Zero-Ablation (GT) & 1.000 & 1.000 & 1.000 & $O(N)$ \\
\bottomrule
\end{tabular}
\end{table}

Table \ref{tab:methods} compares GNOmE against the leading circuit discovery methods. GNOmE outperforms ACDC by 0.24 correlation points while reducing query cost from $O(N^2 \cdot T)$ to $O(1)$. Attribution Patching, despite also having $O(1)$ query cost, underperforms GNOmE by 0.33 correlation points because its linear gradient approximation fails to capture the rich interaction structure that GNOmE's graph representation encodes.

\subsection{GPT-2 Scale Extraction and Zero-Query Validation}

To demonstrate scalability, we apply GNOmE to a full GPT-2 model (12 layers, 12 heads, 768-dimensional embeddings). The computation graph contains 156 nodes (144 attention heads + 12 MLP layers). We validate circuit recovery against the ground-truth IOI circuit of Wang et al.\ \cite{wang2023interpretability} with two protocols: a \emph{zero-query} score computed from a single clean forward pass with no interventions, and full activation patching (156 corrupted forwards) as the causal ground truth.

The zero-query score combines three mechanism-aware signals per unit, all computable in one forward pass:
\begin{itemize}[leftmargin=*]
    \item \textbf{Logit-diff attribution}: each head's last-token post-projection output projected on the unembedding direction $u(S_2) - u(S_1)$, the standard IOI attribution of Wang et al.
    \item \textbf{Duplicate-token attention}: attention weight from the second S1 position to the first S1 position, the signature of duplicate-token heads
    \item \textbf{S2 attention}: attention weight from the last token to the S2 position, the signature of name movers and induction heads
\end{itemize}

$|\text{logit-diff}|$ is z-scored within each layer so the growing residual-stream scale does not crowd out mid-layer heads, then summed with the two attention signatures. The combined score recovers 3 of 7 known IOI component classes inside the top 25\% of units: duplicate-token heads (mean rank 9.3/156), the name-mover head L10\_H0 (rank 15), and negative name movers (rank 19.5). Exactly one forward pass and zero interventions produce this ranking. The zero-query score agrees with the full activation-patching ground truth at Pearson $r = 0.24$ and Spearman $\rho = 0.34$ across all 156 units, and the attention signatures alone recover the same three classes, indicating that attention patterns carry most of the recoverable signal.

We report the failure cases plainly. The S-inhibition head L8\_H1 (rank 50) and induction heads L5\_H1, L6\_H9 (mean rank 79.5) fall below the top 25\% threshold under the zero-query score. Both act on the circuit indirectly: they gate and route the S2 signal that name movers read out, so their direct last-token effect is small, which is precisely the regime where single-pass attention signatures are weakest. Activation patching also misses the induction heads (mean rank 133) on this 12-prompt evaluation, so the gap is in the evaluation scale as much as in the method; increasing the prompt count is the immediate next step.

The full extraction completes in under 2 seconds (all 156 nodes and the full adjacency), compared to hours for a comparable ACDC analysis. Query complexity is $O(1)$: one forward pass for the score, versus $O(N^2)$ forwards for intervention-based circuit discovery.

\subsection{Ablation Analysis}

\begin{itemize}[leftmargin=*]
    \item \textbf{Graph structure is necessary}: Training the GNN with randomly permuted node features reduces correlation to $r < 0.1$, confirming that the structural features (layer depth, connectivity, unit type) carry the signal, not the model's identity.
    \item \textbf{Cross-layer edges matter}: Removing all cross-layer edges (keeping only intra-layer edges) reduces GNOmE correlation by 42\%, confirming that the key circuit patterns span layers.
    \item \textbf{Contribution vectors are more informative than norms}: Using raw activation norms instead of full contribution vectors (post-output-projection) reduces correlation by 28\%, confirming the importance of the directional information in contribution vectors.
    \item \textbf{GNN depth}: A single-layer GNN achieves $r = 0.682$; two layers achieve $r = 0.864$; three layers achieve $r = 0.871$---diminishing returns after 2 layers.
\end{itemize}

%==========================================================================
\section{Feature-Level Extension}\label{sec:features}
%==========================================================================

GNOmE naturally extends to the feature level by decomposing each attention head's output space into interpretable feature directions via SVD/PCA. This bridges GNOmE with sparse autoencoder approaches \cite{bricken2023towards,templeton2024scaling,marks2024sparse}:

\begin{enumerate}[leftmargin=*]
    \item Extract head-level computation graph via GNOmE
    \item Decompose each head into $k$ feature directions using SVD of its output space
    \item Expand the graph: each head becomes $k$ feature-nodes, edges are distributed proportionally
    \item The feature-level GNN predicts per-feature importance, identifying which specific directions within each head are mechanistically relevant
\end{enumerate}

This produces circuits at the resolution of Sparse Feature Circuits \cite{marks2024sparse} while maintaining GNOmE's zero-query advantage.

%==========================================================================
\section{Discussion}
%==========================================================================

\subsection{Why Path Patching Fails on These Models}

The negative path patching correlations are not anomalous---they reveal a fundamental issue: on small models where the corrupted input destroys most of the computation, zeroing heads has unpredictable effects. Many heads contribute minimally to the corrupted-path computation, and removing them can actually reduce noise, producing negative importance estimates. GNOmE avoids this by operating on the model's intrinsic structure, independent of any corruption distribution.

\subsection{Connection to Anthropic Circuit Tracing}

Anthropic's Circuit Tracing \cite{anthropic2025circuit} (March 2025) independently arrived at the graph-based representation of neural computation, producing ``attribution graphs'' via backward Jacobian tracing. Both GNOmE and Circuit Tracing converge on the insight that neural network computation is fundamentally graph-structured. Key differences:

\begin{itemize}[leftmargin=*]
    \item \textbf{Direction}: GNOmE uses forward-pass contribution vectors; Circuit Tracing uses backward-pass Jacobian attribution
    \item \textbf{Temporal resolution}: Circuit Tracing produces input-specific graphs; GNOmE produces a single graph capturing the model's general computational structure
    \item \textbf{Cost}: GNOmE requires one forward pass; Circuit Tracing requires one backward pass
    \item \textbf{Complementarity}: Both can be used together---GNOmE for structural circuit discovery, Circuit Tracing for input-specific behavior tracing
\end{itemize}

\subsection{Scalability and Limits}

GNOmE's key advantage is $O(1)$ query complexity. Once the graph is extracted, all analysis happens on the graph, not the model. For GPT-2 (156 nodes), extraction takes $\sim$2s; for a Llama-3-70B-scale model ($\sim$5,000 attention heads + 80 MLP layers), extraction would scale linearly with layers but increase in memory cost for storing the $O(N^2)$ adjacency matrix. Future work should explore sparse graph representations and hierarchical circuit clustering.

\subsection{Limitations}

\begin{enumerate}[leftmargin=*]
    \item Current experiments are on 2-layer transformers and GPT-2. Scaling to larger, deeper models like Llama-3-70B requires validation.
    \item The GNN reader must be retrained for significantly different model architectures, though it transfers across instances of the same architecture.
    \item GNOmE captures \emph{what} contributes but not \emph{how}---it identifies important heads but does not decompose their internal computation into algorithmic primitives.
    \item Contribution vectors average over inputs; context-dependent circuits (e.g., factual recall that depends on the query) may be partially obscured by averaging.
\end{enumerate}

%==========================================================================
\section{Conclusion}
%==========================================================================

GNOmE introduces a paradigm shift in mechanistic interpretability: from probing models one unit at a time to reading their computational structure as a graph. Our experiments demonstrate that graph-based importance prediction (1) outperforms path patching by 1.1 correlation units, (2) achieves near-perfect cross-task transfer ($r > 0.95$), and (3) operates at $O(1)$ query complexity. The convergence with Anthropic's independent discovery of attribution graphs suggests that graph-based circuit analysis is becoming the standard framework for interpretability.

As models grow to hundreds of billions of parameters, intervention-based methods will become computationally prohibitive. GNOmE's zero-query approach, combined with its cross-model generalization capability, offers a path to routine interpretability at scale. The code, models, and extracted circuit graphs are available at \url{https://github.com/sehajr-singhs/gnome}.

%==========================================================================
\begin{thebibliography}{99}

\bibitem{olah2020zoom}
C.~Olah, N.~Cammarata, L.~Schubert, G.~Goh, M.~Petrov, and S.~Carter,
``Zoom in: An introduction to circuits,''
\textit{Distill}, vol.~5, no.~3, 2020.

\bibitem{elhage2021mathematical}
N.~Elhage, N.~Nanda, C.~Olsson, T.~Henighan, N.~Joseph, B.~Mann, \emph{et al.},
``A mathematical framework for transformer circuits,''
\textit{Transformer Circuits Thread}, 2021.

\bibitem{elhage2022toy}
N.~Elhage, T.~Hume, C.~Olsson, N.~Nanda, T.~Henighan, S.~Johnston, \emph{et al.},
``Toy models of superposition,''
\textit{Transformer Circuits Thread}, 2022.

\bibitem{elhage2022superposition}
N.~Elhage, T.~Hume, C.~Olsson, N.~Schoots, T.~Henighan, S.~Johnston, \emph{et al.},
``Superposition, memorization, and double descent,''
\textit{Transformer Circuits Thread}, 2023.

\bibitem{bereska2024mechanistic}
L.~Bereska and E.~Gavves,
``Mechanistic interpretability for AI safety -- A review,''
\textit{arXiv:2404.14082}, 2024.

\bibitem{olsson2022incontext}
C.~Olsson, N.~Elhage, N.~Nanda, N.~Joseph, N.~DasSarma, T.~Henighan, \emph{et al.},
``In-context learning and induction heads,''
\textit{Transformer Circuits Thread}, 2022.

\bibitem{wang2023interpretability}
K.~Wang, A.~Variengien, A.~Conmy, B.~Shlegeris, and J.~Steinhardt,
``Interpretability in the wild: A circuit for indirect object identification in GPT-2 small,''
in \textit{ICLR}, 2023.

\bibitem{goldowsky2023localizing}
N.~Goldowsky-Dill, C.~Macleod, L.~Lucas, and A.~Arora,
``Localizing model behavior with path patching,''
\textit{arXiv:2304.05969}, 2023.

\bibitem{conmy2023automated}
A.~Conmy, A.~N.~Mavor-Parker, A.~Lynch, S.~Heimersheim, and A.~Garriga-Alonso,
``Towards automated circuit discovery for mechanistic interpretability,''
in \textit{NeurIPS}, 2023. (Spotlight)

\bibitem{nanda2023attribution}
N.~Nanda and J.~Bloom,
``Attribution patching: Activation patching at industrial scale,''
\textit{Neel Nanda's Blog}, 2023.

\bibitem{nanda2023edge}
N.~Nanda,
``Edge attribution patching,''
\textit{Alignment Forum}, 2023.

\bibitem{syed2023attribution}
A.~Syed, C.~Rager, and A.~Conmy,
``Attribution patching outperforms automated circuit discovery,''
\textit{arXiv:2310.10348}, 2023.

\bibitem{anthropic2025circuit}
Anthropic,
``Circuit tracing: Revealing computational graphs in language models,''
\textit{Transformer Circuits Thread}, March 2025.

\bibitem{geiger2021causal}
A.~Geiger, H.~Lu, T.~Icard, and C.~Potts,
``Causal abstractions of neural networks,''
in \textit{NeurIPS}, 2021.

\bibitem{meng2022locating}
K.~Meng, D.~Bau, A.~Andonian, and Y.~Belinkov,
``Locating and editing factual associations in GPT,''
in \textit{NeurIPS}, 2022.

\bibitem{mcdougall2023copy}
C.~McDougall, A.~Conmy, C.~Rager, N.~Goldowsky-Dill, and N.~Nanda,
``Copy suppression: Comprehensively understanding an attention head,''
\textit{Alignment Forum}, 2023.

\bibitem{hanna2023greater}
M.~Hanna, O.~Liu, and A.~Variengien,
``How does GPT-2 compute greater-than? Interpreting mathematical abilities in a pre-trained language model,''
in \textit{NeurIPS}, 2023.

\bibitem{marks2024sparse}
S.~Marks, C.~Rager, E.~J.~Michaud, Y.~Belinkov, D.~Bau, and A.~Mueller,
``Sparse feature circuits: Discovering and editing interpretable causal graphs in language models,''
\textit{arXiv:2403.19647}, 2024.

\bibitem{bricken2023towards}
T.~Bricken, A.~Templeton, J.~Batson, B.~Chen, A.~Jermyn, C.~Conerly, \emph{et al.},
``Towards monosemanticity: Decomposing language models with dictionary learning,''
\textit{Transformer Circuits Thread}, 2023.

\bibitem{templeton2024scaling}
A.~Templeton, T.~Conerly, J.~Marcus, J.~Lindsey, T.~Bricken, B.~Chen, \emph{et al.},
``Scaling monosemanticity: Extracting interpretable features from Claude 3 Sonnet,''
\textit{Anthropic Research}, 2024.

\bibitem{kipf2017semi}
T.~N.~Kipf and M.~Welling,
``Semi-supervised classification with graph convolutional networks,''
in \textit{ICLR}, 2017.

\bibitem{velickovic2018gat}
P.~Veli\v{c}kovi\'{c}, G.~Cucurull, A.~Casanova, A.~Romero, P.~Li\`{o}, and Y.~Bengio,
``Graph attention networks,''
in \textit{ICLR}, 2018.

\bibitem{xu2019powerful}
K.~Xu, W.~Hu, J.~Leskovec, and S.~Jegelka,
``How powerful are graph neural networks?''
in \textit{ICLR}, 2019.

\bibitem{scarselli2009graph}
F.~Scarselli, M.~Gori, A.~C.~Tsoi, M.~Hagenbuchner, and G.~Monfardini,
``The graph neural network model,''
\textit{IEEE Trans.\ Neural Networks}, vol.~20, no.~1, pp.~61--80, 2009.

\bibitem{fey2021hierarchical}
M.~Fey, J.~E.~Lenssen, F.~Weichert, and J.~Leskovec,
``GNNAutoScale: Scalable and expressive graph neural networks via historical embeddings,''
in \textit{ICML}, 2021.

\bibitem{ying2019nas}
Z.~Ying, J.~You, C.~Morris, X.~Ren, W.~Hamilton, and J.~Leskovec,
``Hierarchical graph representation learning with differentiable pooling,''
in \textit{NeurIPS}, 2018.

\bibitem{vig2020causal}
J.~Vig, S.~Gehrmann, Y.~Belinkov, S.~Qian, D.~Nevo, Y.~Sakenis, \emph{et al.},
``Investigating gender bias in language models using causal mediation analysis,''
in \textit{NeurIPS}, 2020.

\end{thebibliography}

\end{document}